\section*{Golden Rules}
Below are some tips based on reflective comments from previous module group projects. Read them carefully; they are crucial for a successful project.
\begin{itemize}
    \item \textbf{Work consistently throughout the project.} Start working on your project promptly and avoid leaving the bulk of the work until just before the deadlines. This group project module is worth 3 credits, and according to official University Guidelines, each credit corresponds to approximately 120 hours of study time. Therefore, you should expect to dedicate around 80 hours in total to your project, averaging about 5 hours per week during term time. If you are not spending this amount of time weekly, you are likely not working sufficiently and may encounter difficulties later.
    \item \textbf{Ensure you have enough to contribute.} Do not wait to be assigned tasks; instead, volunteer for activities that interest you. Remember that your peers' assessment of your contributions significantly impacts your final grade. Therefore, actively participate in the project to be seen as a valuable group member. Take initiative, be cooperative, prioritize the group's and project's interests, and work diligently.
    \item \textbf{Utilize appropriate tools to maximize efficiency.} In particular, implement a version control system to avoid wasting time and effort.
    \item \textbf{Do not underestimate the time required for writing up your report.} On average, written reports constitute 60\% of the total module grade. Allocate sufficient time for this crucial aspect.
    \item \textbf{Monitor the progress of all group members.} Keep track of each member's progress and ensure everyone is consistently working on their assigned tasks. Ultimately, it is the group's responsibility to monitor individual progress, ensure tasks are completed as required, and take corrective action when necessary.
    \item \textbf{Take meetings seriously.} Regular formal and informal meetings are essential for project success. Dedicate time to prepare for each meeting and take personal notes during the session in addition to the official minutes. However, note that the majority of the actual project work will be conducted outside of meetings. Students who only contribute during meetings will find they have limited content for their reports, which will likely result in lower module grades.
    \item \textbf{Maintain thorough records of your meetings.} Compiling your minutes on a shared platform, such as a webpage, is an effective way to prevent loss and ensure accessibility for all members. Use the minutes to ensure that the group avoids repetitive discussions, remembers decisions made, and maintains consistent progress.
\end{itemize}